\section{Experiment}

We tested the created courses on a sample of computer science students at \ac{fiit}. The students were all in their second year and attending a software modeling class. Each student was given an account on our Moodle site and was enrolled in one of the courses. We decided to assign the course randomly. It is likely that if we let students choose the course variant on their own, only those who enjoy competitive activities would choose the competitive variant. That might give us biased results with overly positive reception of the competitive approach to educational content - which has been demonstrated to be a potentially polarizing factor, as not some people find it more unpleasant than motivating\todo{ref if possible}.

Participation in the experiment was voluntary. Students who filled in their assigned course were given bonus points (2) to their software modeling class, and those who scored in top 30\% across both courses were given double the credits (4). In total, 35 students completed their course - 14 for the competitive variant, 21 for the playful variant. Some of the participants dropped out, resulting in varying amounts of activity completions throughout the courses. 

We collected data describing the participants' behavior, objective performance and subjective perception. The courses were available to students between 7.4.2026 and 24.4.2026.